% ============================================================
% fig_pipeline_minimo.tex: Pipeline del Transformer mínimo
% Uso: % ============================================================
% fig_pipeline_minimo.tex: Pipeline del Transformer mínimo
% Uso: % ============================================================
% fig_pipeline_minimo.tex: Pipeline del Transformer mínimo
% Uso: % ============================================================
% fig_pipeline_minimo.tex: Pipeline del Transformer mínimo
% Uso: \input{figures/fig_pipeline_minimo}
% ============================================================
\begin{figure}[H]
\centering
\begin{tikzpicture}[
    blk/.style={draw, rounded corners, minimum width=2.4cm,
                minimum height=0.7cm, font=\small, align=center},
    inp/.style={blk, fill=gray!10,    draw=gray!50},
    emb/.style={blk, fill=purple!10,  draw=purple!40},
    posblk/.style={blk, fill=teal!10,    draw=teal!50},
    enc/.style={blk, fill=blue!10,    draw=blue!40},
    outblk/.style={blk, fill=orange!10,  draw=orange!50},
    sft/.style={blk, fill=green!10,   draw=green!50!black},
    arrow/.style={-{Stealth}, thick},
    lbl/.style={font=\tiny, gray, align=center}
]

% ── Tokens de entrada ────────────────────────────────────────
\node[inp] (chars) at (0, 0) {``hola''\\{\tiny $[h,o,l,a]$}};
\node[lbl] at (0, -0.75) {caracteres};

% ── Índices enteros ──────────────────────────────────────────
\node[inp] (idx) at (2.8, 0) {índices\\{\tiny $[7,14,11,0]$}};
\node[lbl] at (2.8, -0.75) {vocabulario $|\mathcal{V}|$};

% ── Embedding ─────────────────────────────────────────────────
\node[emb] (emb) at (5.6, 0) {Embedding\\{\tiny $\mathbb{R}^{n \times d_{\text{model}}}$}};
\node[lbl] at (5.6, -0.75) {$E \in \mathbb{R}^{|\mathcal{V}| \times d}$};

% ── Positional Encoding ───────────────────────────────────────
\node[posblk] (pe) at (8.4, 0) {Positional\\Encoding};
\node[lbl] at (8.4, -0.75) {$\bm{x}'_t = \bm{e}_t + PE_t$};

% ── Encoder (N bloques) ───────────────────────────────────────
\node[enc] (enc) at (11.2, 0) {Encoder\\{\tiny $\times N$ bloques}};
\node[lbl] at (11.2, -0.75) {MHA + FFN + Norm};

% ── Proyección lineal ─────────────────────────────────────────
\node[outblk] (lin) at (11.2, -2.0) {Linear\\{\tiny $\mathbb{R}^{d} \to \mathbb{R}^{|\mathcal{V}|}$}};

% ── Softmax ──────────────────────────────────────────────────
\node[sft] (sft) at (8.4, -2.0) {Softmax};
\node[lbl] at (8.4, -2.75) {$P(\text{siguiente car.})$};

% ── Token predicho ────────────────────────────────────────────
\node[inp] (pred) at (5.6, -2.0) {``?''\\{\tiny token predicho}};
\node[lbl] at (5.6, -2.75) {argmax o muestreo};

% ── Flechas forward ───────────────────────────────────────────
\draw[arrow] (chars) -- (idx);
\draw[arrow] (idx)   -- (emb);
\draw[arrow] (emb)   -- (pe);
\draw[arrow] (pe)    -- (enc);
\draw[arrow] (enc)   -- (lin);
\draw[arrow] (lin)   -- (sft);
\draw[arrow] (sft)   -- (pred);

% ── Anotación: máscara causal ─────────────────────────────────
\node[draw=blue!20, fill=blue!4, rounded corners, font=\tiny,
      inner sep=4pt, align=left] at (11.2, 1.35)
    {M\'ascara causal:\\el token $t$ solo\\ve $1,\ldots,t{-}1$};
\draw[-{Stealth}, gray!40, thin] (11.2, 0.65) -- (11.2, 0.95);

% ── Objetivo de entrenamiento ─────────────────────────────────
\node[draw=orange!30, fill=orange!4, rounded corners, font=\tiny,
      inner sep=5pt, align=center] at (3.0, -2.0)
    {Entrenamiento:\\$\mathcal{L} = -\log P(c_{t+1} \mid c_1,\ldots,c_t)$\\
     Entrop\'ia cruzada sobre el\\siguiente car\'acter correcto};

% ── Flechas de flujo de la pérdida ────────────────────────────
\draw[-{Stealth}, orange!50, dashed, thick]
    (pred.west) -- +(-0.3, 0) -- +(-0.3, -0.7)
    node[below, font=\tiny, orange!70] {$\hat{y}$ vs $y$};

\end{tikzpicture}
\caption{Pipeline completo del Transformer m\'inimo implementado en este
trabajo. La secuencia de caracteres de entrada se convierte en \'indices
enteros, se proyecta al espacio de embeddings, se le suma la codificaci\'on
posicional, y se procesa a trav\'es de $N$ bloques encoder con m\'ascara
causal. La proyecci\'on lineal final y el softmax producen una distribuci\'on
de probabilidad sobre el vocabulario de caracteres. Durante el entrenamiento,
la p\'erdida de entrop\'ia cruzada compara la predicci\'on con el siguiente
car\'acter real de la secuencia.}
\label{fig:pipeline_minimo}
\end{figure}

% ============================================================
\begin{figure}[H]
\centering
\begin{tikzpicture}[
    blk/.style={draw, rounded corners, minimum width=2.4cm,
                minimum height=0.7cm, font=\small, align=center},
    inp/.style={blk, fill=gray!10,    draw=gray!50},
    emb/.style={blk, fill=purple!10,  draw=purple!40},
    posblk/.style={blk, fill=teal!10,    draw=teal!50},
    enc/.style={blk, fill=blue!10,    draw=blue!40},
    outblk/.style={blk, fill=orange!10,  draw=orange!50},
    sft/.style={blk, fill=green!10,   draw=green!50!black},
    arrow/.style={-{Stealth}, thick},
    lbl/.style={font=\tiny, gray, align=center}
]

% ── Tokens de entrada ────────────────────────────────────────
\node[inp] (chars) at (0, 0) {``hola''\\{\tiny $[h,o,l,a]$}};
\node[lbl] at (0, -0.75) {caracteres};

% ── Índices enteros ──────────────────────────────────────────
\node[inp] (idx) at (2.8, 0) {índices\\{\tiny $[7,14,11,0]$}};
\node[lbl] at (2.8, -0.75) {vocabulario $|\mathcal{V}|$};

% ── Embedding ─────────────────────────────────────────────────
\node[emb] (emb) at (5.6, 0) {Embedding\\{\tiny $\mathbb{R}^{n \times d_{\text{model}}}$}};
\node[lbl] at (5.6, -0.75) {$E \in \mathbb{R}^{|\mathcal{V}| \times d}$};

% ── Positional Encoding ───────────────────────────────────────
\node[posblk] (pe) at (8.4, 0) {Positional\\Encoding};
\node[lbl] at (8.4, -0.75) {$\bm{x}'_t = \bm{e}_t + PE_t$};

% ── Encoder (N bloques) ───────────────────────────────────────
\node[enc] (enc) at (11.2, 0) {Encoder\\{\tiny $\times N$ bloques}};
\node[lbl] at (11.2, -0.75) {MHA + FFN + Norm};

% ── Proyección lineal ─────────────────────────────────────────
\node[outblk] (lin) at (11.2, -2.0) {Linear\\{\tiny $\mathbb{R}^{d} \to \mathbb{R}^{|\mathcal{V}|}$}};

% ── Softmax ──────────────────────────────────────────────────
\node[sft] (sft) at (8.4, -2.0) {Softmax};
\node[lbl] at (8.4, -2.75) {$P(\text{siguiente car.})$};

% ── Token predicho ────────────────────────────────────────────
\node[inp] (pred) at (5.6, -2.0) {``?''\\{\tiny token predicho}};
\node[lbl] at (5.6, -2.75) {argmax o muestreo};

% ── Flechas forward ───────────────────────────────────────────
\draw[arrow] (chars) -- (idx);
\draw[arrow] (idx)   -- (emb);
\draw[arrow] (emb)   -- (pe);
\draw[arrow] (pe)    -- (enc);
\draw[arrow] (enc)   -- (lin);
\draw[arrow] (lin)   -- (sft);
\draw[arrow] (sft)   -- (pred);

% ── Anotación: máscara causal ─────────────────────────────────
\node[draw=blue!20, fill=blue!4, rounded corners, font=\tiny,
      inner sep=4pt, align=left] at (11.2, 1.35)
    {M\'ascara causal:\\el token $t$ solo\\ve $1,\ldots,t{-}1$};
\draw[-{Stealth}, gray!40, thin] (11.2, 0.65) -- (11.2, 0.95);

% ── Objetivo de entrenamiento ─────────────────────────────────
\node[draw=orange!30, fill=orange!4, rounded corners, font=\tiny,
      inner sep=5pt, align=center] at (3.0, -2.0)
    {Entrenamiento:\\$\mathcal{L} = -\log P(c_{t+1} \mid c_1,\ldots,c_t)$\\
     Entrop\'ia cruzada sobre el\\siguiente car\'acter correcto};

% ── Flechas de flujo de la pérdida ────────────────────────────
\draw[-{Stealth}, orange!50, dashed, thick]
    (pred.west) -- +(-0.3, 0) -- +(-0.3, -0.7)
    node[below, font=\tiny, orange!70] {$\hat{y}$ vs $y$};

\end{tikzpicture}
\caption{Pipeline completo del Transformer m\'inimo implementado en este
trabajo. La secuencia de caracteres de entrada se convierte en \'indices
enteros, se proyecta al espacio de embeddings, se le suma la codificaci\'on
posicional, y se procesa a trav\'es de $N$ bloques encoder con m\'ascara
causal. La proyecci\'on lineal final y el softmax producen una distribuci\'on
de probabilidad sobre el vocabulario de caracteres. Durante el entrenamiento,
la p\'erdida de entrop\'ia cruzada compara la predicci\'on con el siguiente
car\'acter real de la secuencia.}
\label{fig:pipeline_minimo}
\end{figure}

% ============================================================
\begin{figure}[H]
\centering
\begin{tikzpicture}[
    blk/.style={draw, rounded corners, minimum width=2.4cm,
                minimum height=0.7cm, font=\small, align=center},
    inp/.style={blk, fill=gray!10,    draw=gray!50},
    emb/.style={blk, fill=purple!10,  draw=purple!40},
    posblk/.style={blk, fill=teal!10,    draw=teal!50},
    enc/.style={blk, fill=blue!10,    draw=blue!40},
    outblk/.style={blk, fill=orange!10,  draw=orange!50},
    sft/.style={blk, fill=green!10,   draw=green!50!black},
    arrow/.style={-{Stealth}, thick},
    lbl/.style={font=\tiny, gray, align=center}
]

% ── Tokens de entrada ────────────────────────────────────────
\node[inp] (chars) at (0, 0) {``hola''\\{\tiny $[h,o,l,a]$}};
\node[lbl] at (0, -0.75) {caracteres};

% ── Índices enteros ──────────────────────────────────────────
\node[inp] (idx) at (2.8, 0) {índices\\{\tiny $[7,14,11,0]$}};
\node[lbl] at (2.8, -0.75) {vocabulario $|\mathcal{V}|$};

% ── Embedding ─────────────────────────────────────────────────
\node[emb] (emb) at (5.6, 0) {Embedding\\{\tiny $\mathbb{R}^{n \times d_{\text{model}}}$}};
\node[lbl] at (5.6, -0.75) {$E \in \mathbb{R}^{|\mathcal{V}| \times d}$};

% ── Positional Encoding ───────────────────────────────────────
\node[posblk] (pe) at (8.4, 0) {Positional\\Encoding};
\node[lbl] at (8.4, -0.75) {$\bm{x}'_t = \bm{e}_t + PE_t$};

% ── Encoder (N bloques) ───────────────────────────────────────
\node[enc] (enc) at (11.2, 0) {Encoder\\{\tiny $\times N$ bloques}};
\node[lbl] at (11.2, -0.75) {MHA + FFN + Norm};

% ── Proyección lineal ─────────────────────────────────────────
\node[outblk] (lin) at (11.2, -2.0) {Linear\\{\tiny $\mathbb{R}^{d} \to \mathbb{R}^{|\mathcal{V}|}$}};

% ── Softmax ──────────────────────────────────────────────────
\node[sft] (sft) at (8.4, -2.0) {Softmax};
\node[lbl] at (8.4, -2.75) {$P(\text{siguiente car.})$};

% ── Token predicho ────────────────────────────────────────────
\node[inp] (pred) at (5.6, -2.0) {``?''\\{\tiny token predicho}};
\node[lbl] at (5.6, -2.75) {argmax o muestreo};

% ── Flechas forward ───────────────────────────────────────────
\draw[arrow] (chars) -- (idx);
\draw[arrow] (idx)   -- (emb);
\draw[arrow] (emb)   -- (pe);
\draw[arrow] (pe)    -- (enc);
\draw[arrow] (enc)   -- (lin);
\draw[arrow] (lin)   -- (sft);
\draw[arrow] (sft)   -- (pred);

% ── Anotación: máscara causal ─────────────────────────────────
\node[draw=blue!20, fill=blue!4, rounded corners, font=\tiny,
      inner sep=4pt, align=left] at (11.2, 1.35)
    {M\'ascara causal:\\el token $t$ solo\\ve $1,\ldots,t{-}1$};
\draw[-{Stealth}, gray!40, thin] (11.2, 0.65) -- (11.2, 0.95);

% ── Objetivo de entrenamiento ─────────────────────────────────
\node[draw=orange!30, fill=orange!4, rounded corners, font=\tiny,
      inner sep=5pt, align=center] at (3.0, -2.0)
    {Entrenamiento:\\$\mathcal{L} = -\log P(c_{t+1} \mid c_1,\ldots,c_t)$\\
     Entrop\'ia cruzada sobre el\\siguiente car\'acter correcto};

% ── Flechas de flujo de la pérdida ────────────────────────────
\draw[-{Stealth}, orange!50, dashed, thick]
    (pred.west) -- +(-0.3, 0) -- +(-0.3, -0.7)
    node[below, font=\tiny, orange!70] {$\hat{y}$ vs $y$};

\end{tikzpicture}
\caption{Pipeline completo del Transformer m\'inimo implementado en este
trabajo. La secuencia de caracteres de entrada se convierte en \'indices
enteros, se proyecta al espacio de embeddings, se le suma la codificaci\'on
posicional, y se procesa a trav\'es de $N$ bloques encoder con m\'ascara
causal. La proyecci\'on lineal final y el softmax producen una distribuci\'on
de probabilidad sobre el vocabulario de caracteres. Durante el entrenamiento,
la p\'erdida de entrop\'ia cruzada compara la predicci\'on con el siguiente
car\'acter real de la secuencia.}
\label{fig:pipeline_minimo}
\end{figure}

% ============================================================
\begin{figure}[H]
\centering
\begin{tikzpicture}[
    blk/.style={draw, rounded corners, minimum width=2.4cm,
                minimum height=0.7cm, font=\small, align=center},
    inp/.style={blk, fill=gray!10,    draw=gray!50},
    emb/.style={blk, fill=purple!10,  draw=purple!40},
    posblk/.style={blk, fill=teal!10,    draw=teal!50},
    enc/.style={blk, fill=blue!10,    draw=blue!40},
    outblk/.style={blk, fill=orange!10,  draw=orange!50},
    sft/.style={blk, fill=green!10,   draw=green!50!black},
    arrow/.style={-{Stealth}, thick},
    lbl/.style={font=\tiny, gray, align=center}
]

% ── Tokens de entrada ────────────────────────────────────────
\node[inp] (chars) at (0, 0) {``hola''\\{\tiny $[h,o,l,a]$}};
\node[lbl] at (0, -0.75) {caracteres};

% ── Índices enteros ──────────────────────────────────────────
\node[inp] (idx) at (2.8, 0) {índices\\{\tiny $[7,14,11,0]$}};
\node[lbl] at (2.8, -0.75) {vocabulario $|\mathcal{V}|$};

% ── Embedding ─────────────────────────────────────────────────
\node[emb] (emb) at (5.6, 0) {Embedding\\{\tiny $\mathbb{R}^{n \times d_{\text{model}}}$}};
\node[lbl] at (5.6, -0.75) {$E \in \mathbb{R}^{|\mathcal{V}| \times d}$};

% ── Positional Encoding ───────────────────────────────────────
\node[posblk] (pe) at (8.4, 0) {Positional\\Encoding};
\node[lbl] at (8.4, -0.75) {$\bm{x}'_t = \bm{e}_t + PE_t$};

% ── Encoder (N bloques) ───────────────────────────────────────
\node[enc] (enc) at (11.2, 0) {Encoder\\{\tiny $\times N$ bloques}};
\node[lbl] at (11.2, -0.75) {MHA + FFN + Norm};

% ── Proyección lineal ─────────────────────────────────────────
\node[outblk] (lin) at (11.2, -2.0) {Linear\\{\tiny $\mathbb{R}^{d} \to \mathbb{R}^{|\mathcal{V}|}$}};

% ── Softmax ──────────────────────────────────────────────────
\node[sft] (sft) at (8.4, -2.0) {Softmax};
\node[lbl] at (8.4, -2.75) {$P(\text{siguiente car.})$};

% ── Token predicho ────────────────────────────────────────────
\node[inp] (pred) at (5.6, -2.0) {``?''\\{\tiny token predicho}};
\node[lbl] at (5.6, -2.75) {argmax o muestreo};

% ── Flechas forward ───────────────────────────────────────────
\draw[arrow] (chars) -- (idx);
\draw[arrow] (idx)   -- (emb);
\draw[arrow] (emb)   -- (pe);
\draw[arrow] (pe)    -- (enc);
\draw[arrow] (enc)   -- (lin);
\draw[arrow] (lin)   -- (sft);
\draw[arrow] (sft)   -- (pred);

% ── Anotación: máscara causal ─────────────────────────────────
\node[draw=blue!20, fill=blue!4, rounded corners, font=\tiny,
      inner sep=4pt, align=left] at (11.2, 1.35)
    {M\'ascara causal:\\el token $t$ solo\\ve $1,\ldots,t{-}1$};
\draw[-{Stealth}, gray!40, thin] (11.2, 0.65) -- (11.2, 0.95);

% ── Objetivo de entrenamiento ─────────────────────────────────
\node[draw=orange!30, fill=orange!4, rounded corners, font=\tiny,
      inner sep=5pt, align=center] at (3.0, -2.0)
    {Entrenamiento:\\$\mathcal{L} = -\log P(c_{t+1} \mid c_1,\ldots,c_t)$\\
     Entrop\'ia cruzada sobre el\\siguiente car\'acter correcto};

% ── Flechas de flujo de la pérdida ────────────────────────────
\draw[-{Stealth}, orange!50, dashed, thick]
    (pred.west) -- +(-0.3, 0) -- +(-0.3, -0.7)
    node[below, font=\tiny, orange!70] {$\hat{y}$ vs $y$};

\end{tikzpicture}
\caption{Pipeline completo del Transformer m\'inimo implementado en este
trabajo. La secuencia de caracteres de entrada se convierte en \'indices
enteros, se proyecta al espacio de embeddings, se le suma la codificaci\'on
posicional, y se procesa a trav\'es de $N$ bloques encoder con m\'ascara
causal. La proyecci\'on lineal final y el softmax producen una distribuci\'on
de probabilidad sobre el vocabulario de caracteres. Durante el entrenamiento,
la p\'erdida de entrop\'ia cruzada compara la predicci\'on con el siguiente
car\'acter real de la secuencia.}
\label{fig:pipeline_minimo}
\end{figure}
